% !TEX TS-program = lualatex
% !TEX encoding = UTF-8

\documentclass[Session2026.tex]{subfiles}

\ifcsname preamble@file\endcsname
  \setcounter{page}{\getpagerefnumber{M-04_mercredi}}
\fi

\begin{document}

\addcontentsline{toc}{chapter}{Mercredi 8 juillet, de la férie}

\bigtitle{Mercredi de la 14\textsuperscript{e} semaine \emph{per annum},\\à la Messe}{mercredi 8 juillet}{Messe}

\gscore{in_suscepimus}
\translation{\aa Nous avons reçu, ô Dieu, ta miséricorde au milieu de ton temple ; comme ton nom, ainsi ta louange retentit jusqu'aux confins de la terre ; ta droite est pleine de justice.\\
\vv Grand est le Seigneur, et très digne de louange, dans la cité de notre Dieu, sur sa montagne sainte.}

\smalltitle{Psalmodie de Tierce}

\gscore{an_clamavi}
\translation{\aa J'ai crié, et il m'a exaucé.}
\smalltitle{Psaume 119}
\psalm{119}{E}
\smalltitle{Psaume 120}
\psalm{120}{E}
\needspace{3cm}
\smalltitle{Psaume 121}
\psalm{121}{E}

\smalltitle{Kyrie XVI}
\gscore{ky_k16}

\gscore{gr_esto_mihi}
\translation{\rr Sois pour moi un Dieu protecteur et un lieu de refuge, afin de me sauver.\\
\vv Dieu, j'espère en toi ; Seigneur, je ne serai pas confondu à jamais.}

\gscore{al_magnus_dominus}
\translation{\vv Grand est le Seigneur, et très digne de louange, dans la cité de notre Dieu, sur sa montagne sainte.}
\gscore{of_populum_humilem}
\translation{\rr Tu sauveras le peuple humble, Seigneur, et tu humilieras les yeux des superbes ; car qui est Dieu sinon toi, Seigneur ?}

\smalltitle{Sanctus XVI}
\gscore{ky_s16}
\smalltitle{Agnus XVI}
\gscore{ky_a16}
\gscore{co_gustate}
\translation{\aa Goûtez et voyez comme est doux le Seigneur: heureux l'homme qui espère en lui.\\
\vv Je bénirai le Seigneur en tous temps, sa louange sans cesse à ma bouche.\\
\vv Je me glorifierai dans le Seigneur, que les pauvres m'entendent et se réjouissent.}


\bigtitle{Mercredi II, à Sexte}{mercredi 8 juillet}{Sexte}

\rubric{Ouverture, Hymne et Psalmodie de Sexte, comme au mardi, p. \pageref{hy_rector_potens_ferialis}.}

\capitulum{Col 3, 17}
{Omne quodcúmque fácitis in verbo aut in ópere, † ómnia in nómine
Dómini Iesu * grátias agéntes Deo Patri per ipsum.}
{Tout ce que vous dites, tout ce que vous faites, que ce soit toujours au
nom du Seigneur Jésus-Christ, en offrant par lui votre action de grâce à
Dieu le Père.}

\versiculus{Tibi, Dómine, sacrificábo hóstiam laudis}{Et nomen Dómini invocábo}{Je t’offrirai, Seigneur, le sacrifice d’action de grâce}{J’invoquerai le nom du Seigneur}

\oratio{TODO}{TODO}

\gscore{or_benedicamus_hm_ferialis}

\bigtitle{Mercredi II, à None}{mercredi 8 juillet}{None}

\smallscore{or_dia_ferialis}

\rubric{Ouverture, Hymne et Psalmodie de None, comme au mardi, p. \pageref{an_beati_omnes}.}

\capitulum{Col 3, 23-24}
{Quodcúmque fácitis, ex ánimo operámini sicut Dómino et non
homínibus, * sciéntes quod a Dómino accipiétis retributiónem hereditátis. / Dómino Christo servíte.}
{Quel que soit votre travail, faites-le de bon coeur, pour le Seigneur et non
pour plaire à des hommes : vous savez bien qu’en retour le Seigneur fera
de vous ses héritiers. Le maître c’est le Christ : vous êtes à son service.}

\versiculus{Dóminus pars hereditátis meæ et cálicis mei}{Tu es qui détines sortem meam}{Seigneur, mon partage et ma coupe}{De toi dépend mon sort}

\rubric{Oraison des Vêpres, p. \pageref{0709_oratio}.}

\gscore{or_benedicamus_hm_ferialis}

\bigtitle{Vierge Marie, Mère de Providence,\\aux Vêpres}{mercredi 9 juillet}{Vêpres}

\smallscore{or_dia_ferialis}

\smalltitle{Psaume 134}
\gscore{an_omnia_quaecumque}
\translation{\aa Tout ce qu'il veut, le Seigneur le fait.}
\psalm{134}{3}

\smalltitle{Psaume 135}
\gscore{an_quoniam_in_aeternum}
\translation{\aa Car éternel est son amour.}
\psalm{135}{3}

\smalltitle{Psaume 136}
\gscore{an_hymnum_cantate}
\translation{\aa Chantez-nous un chant de Sion.}
\psalm{136}{8}

\smalltitle{Psaume 137}
\gscore{an_in_conspectu_angelorum}
\translation{\aa Mon dieu, je te chanterai des psaumes en présence des anges.}
\psalm{137}{5}

\capitulum{1 Jn 2, 3-6}
{In hoc cognóscimus quóniam nóvimus Christum : * si mandáta eius
servémus. / Qui dicit : « Novi eum », et mandáta eius non servat, mendax
est, et in isto véritas non est ; † qui autem servat verbum eius, * vere in
hoc cáritas Dei consummáta est. / In hoc cognóscimus quóniam in ipso
sumus. / Qui dicit se in ipso manére, debet, sicut ille ambulávit, et ipse
ambuláre.}
{Voici comment nous pouvons savoir que nous connaissons Jésus-Christ :
c’est en gardant ses commandements. Celui qui dit : « Je le connais »,
et qui ne garde pas ses commandements, est un menteur : la vérité n’est
pas en lui. Mais en celui qui garde fidèlement sa parole, l’amour de Dieu
atteint vraiment la perfection : voilà comment nous reconnaissons que
nous sommes en lui. Celui qui déclare demeurer en lui doit marcher lui-même
dans la voie où lui, Jésus, a marché.}

\gscore{rb_custodi_nos}
\translation{\rr Garde-nous, Seigneur, comme la prunelle de l'œil.\\
\vv Protège-nous à l'ombre de tes ailes.}

\gscore{hy_o_quam_glorifica}
\translation{\colored{D}e quelle glorieuse lumière vous brillez, royale enfant de la lignée de David, Vierge Marie, vous qui siègez très haut au dessus de tous les habitants des cieux.\\
\colored{M}ère qui gardez votre honneur de vierge, vous avez préparé chastement, pour le Seigneur des cieux, dans votre sein consacré, la demeure de votre cœur, d'où naquit le Christ, Dieu incarné.\\
\colored{C}elui qu'adore avec vénération tout l'univers, devant qui à bon droit tout genou fléchit, nous implorons de Lui, par votre secours, les joies de la lumière qui chasse les ténèbres. \\
\colored{F}ais-nous cette largesse, Père de toute lumière, dans l'Esprit Saint, par Ton propre Fils qui, vivant avec Toi dans l'éclat de l'éther, règne et gouverne tous les siècles.}

\versiculus{Diffúsa est grátia in lábiis tuis}{Proptérea benedíxit te Deus in ætérnum}{La grâce est répandue sur vos lèvres}{Que Dieu vous bénisse pour toujours}

\smalltitle{Magnificat}
\gscore{an_exsultavit_quia_fecit}
\translation{\aa Mon esprit exulte en Dieu mon sauveur, car le puissant fit pour moi des merveilles.}
\incipit{magn}{6f}
\psalm{magn}{6}

\needspace{3cm}
\smalltitle{Intercession}
\smallscore{or_kyrie_ferialis}
\smallscore{or_pater_ferialis}

\label{0709_oratio}
\oratio{Omnípotens, sempitérne Deus, Ecclésiam tuam, déxtera poténtiæ tuæ, a cunctis prótege perículis,~\pscross et beáta María Matre Providéntiæ intercedénte,~\psstar fac eam præsénti gaudére prosperitáte et futúra. Per Dóminum.}{Dieu éternel et tout puissant, protège ton Église de tous les dangers par la force de ton bras, et par l'intercession de la bienheureuse Marie, Mère de Providence, fais-la se réjouir de sa prospérité présente et future. Par Jésus-Christ.}

\blessing

\gscore{or_benedicamus_festis}

\end{document}
